\section{Discussion}\label{sec:discussion}

\paragraph{An executable gate, not a replacement for review} The oracle and the VLM catch different errors (RQ2). The oracle is deterministic, explains its verdicts and costs nothing per run, but it only sees what the contract exposes. The VLM sees the drawn shapes but cannot see invisible conditions, misreads intermittent motion, and changes its verdict between identical runs. In a generation pipeline we would therefore use the oracle as the \emph{acceptance gate} and repair driver, and keep a visual review as an additional, advisory check for shape errors outside the contract. In V\&V terms, this puts analytic comparison first and animation second \cite{sargent2013vv}. Our data show that a VLM alone as gate is not safe: 12 of its stop decisions accepted simulators that the oracle rejected.

\paragraph{Feedback should carry the reference value when one exists} Iterative repair with test feedback is well established \cite{chen2023selfdebug,xia2024chatrepair,olausson2024selfrepair}, and the quality of feedback is known to limit self-repair \cite{olausson2024selfrepair}. Our ablation isolates one dimension of that quality, \emph{information content}, at fixed source and wording, and shows that its effect depends on the model. Localisation alone sufficed for the larger model, which held the missing design knowledge, but did nothing for the smaller model, which did not. For a mechanism with published dimensions the oracle knows the correct value anyway, so withholding it has no benefit. More generally, an oracle that encodes reference knowledge should return that knowledge in its diagnostics, because the cheapest repairer may need it. This contrasts with feedback designs that report only failing tests or error magnitudes, and with results suggesting that weaker models mainly benefit from more feedback \cite{han2025convcodeworld,pan2025gfr}. In our setting it was the \emph{kind} of information that mattered, not its amount.

\paragraph{Confident wrong knowledge} The most consistent failure we saw was not an inability to edit code. It was a repairer that defended a wrong design with a confident and specific claim (``the lengths match the published holy numbers''), sometimes blaming the checker. Self-correction without external feedback is known to be unreliable \cite{huang2024cannot,kamoi2024selfcorrection}, and LLMs hallucinate domain facts in code \cite{liu2026codehallucination,tian2025codehalu}. Our data add that in an engineering domain such claims are cheap to falsify mechanically, but only if the acceptance step does not depend on the same model's judgement.

\paragraph{When is an oracle worth writing?} The oracle needed one module per mechanism type (27--90 lines of JavaScript each), written once and reused for every generated instance regardless of its dimensions. The checks that need no ground truth (META) already caught most seeded faults (17 of 24), and they transfer to any simulator that follows the contract. Errors were concentrated where design knowledge must be recalled rather than derived. Mechanisms specified by a relation or law that the description implies (cam, Peaucellier--Lipkin) produced no natural errors, whereas mechanisms that depend on remembered numbers (Jansen, lambda) produced most of them. A practical heuristic follows: invest in DESIGN rules for mechanisms whose function depends on published dimensions, and rely on the generic META, PARAM and GT checks elsewhere.

\paragraph{Leakage through magnitudes} Our coarse and location-only arms still reported the measured value of each failed check. The smaller model repeatedly tried to invert this value into a corrected length, without success because the minimax formula was not shown. The larger model succeeded twice on Jansen. Error magnitudes are thus a side channel between feedback levels, and future ablations should control them explicitly.
